% Generated draft: requires mathematical and source review.
\documentclass[11pt]{article}
\usepackage[margin=25mm]{geometry}
\usepackage{fontspec}
\setmainfont{Latin Modern Roman}
\usepackage{amsmath,amssymb,mathtools,unicode-math}
\setmathfont{Latin Modern Math}
\usepackage{enumitem,needspace,xurl,hyperref}
\setlength{\parindent}{0pt}
\setlength{\parskip}{5pt}
\setlength{\emergencystretch}{3em}
\newcommand{\N}{\mathbb{N}}
\newcommand{\Z}{\mathbb{Z}}
\newcommand{\Q}{\mathbb{Q}}
\newcommand{\R}{\mathbb{R}}
\newcommand{\C}{\mathbb{C}}
\newcommand{\F}{\mathbb{F}}
\newcommand{\A}{\mathbb{A}}
\newcommand{\PP}{\mathbb{P}}
\newcommand{\E}{\mathbb{E}}
\newcommand{\eps}{\varepsilon}
\newcommand{\dd}{\,\mathrm{d}}
\begin{document}
\section*{UnsolvedMath ID: 30002298}
\textbf{Source problem code:} catalogue entry
\subsection{Definitions and mathematical statement}
Let \(P\) be a compact polyhedron in the class intended by the source, and write \(V(P)\) for its vertex set. The notation \(\langle v,u\rangle\) denotes the ambient pairing used in the source, with \(u\) the transform variable. Write \(\Omega_P(u)\) for the polynomial called the denominator of the Fantappi\`e transform of \(P\) in the supplied record.

The requested problem is to give a necessary and sufficient geometric characterization of those \(P\) for which
\[
\Omega_P(u)=\prod_{v\in V(P)}\bigl(1-\langle v,u\rangle\bigr).
\]
The equality is a denominator identity in the transform variable, rather than an equality at a single value of \(u\).

The supplied record does not define the Fantappi\`e transform, its underlying measure or density, or the convention selecting its denominator. In particular, it does not specify whether the denominator is reduced after cancellation or how it is normalized. It also does not specify the ambient dimension, whether convexity or full dimensionality is required, or the meaning of vertices if nonconvex polyhedra are allowed. These qualifications remain unresolved; a fully determined characterization problem requires the conventions from the cited report.
\subsection{Short English statement}
Characterize geometrically the compact polyhedra whose Fantappi\`e-transform denominator is exactly the product of the linear factors associated with their vertices. The transform, denominator convention, and intended class of polyhedra require clarification from the source.
\subsection{Sources}
\begin{itemize}
\item Supplied local UnsolvedMath registry record, permanent ID \(30002298\), titled ``Polyhedra with Vertex-Factored Fantappi\`e Denominators''; exact source problem code \texttt{catalogue entry}. The record belongs to the Oberwolfach Reports Open Problems collection and gives proposed year 2013; no individual proposer is identified.
\item Supplied report citation: \emph{Structured Function Systems and Applications} (2013), Oberwolfach Reports, pp.\ 579--655, DOI \texttt{10.4171/owr/2013/11}. Source URL and supplied literature-source link, labeled ``Oberwolfach Report 2013'': \url{https://doi.org/10.4171/owr/2013/11}.
\item The registry labels the statement \texttt{corrected\_verified} and reports that its mathematical content was checked against the official report, with notation, context, and question form normalized. This is supplied verification metadata, not independent verification performed for this draft.
\item Literature assessment dated exactly \texttt{2026-08-21}: the supplied curation labels the problem open and states that no complete characterization is stated or verified in the extracted record. It describes a moment-generating-function context and contextual polyhedral observations, without reproducing those observations. This is supplied dated curation, not an independently verified literature assessment; the record expressly calls for independent verification of time-sensitive status and cited claims.
\end{itemize}
\end{document}
